\begin{lemma}[Fixed-field Hecke Siegel--Walfisz estimate]
\label{lem:gaussian-hecke-sw}
Let $E=\Q(i)$, let $\chi$ be a fixed finite-order ray character of $E$
with finite modulus $\mathfrak m$, and let $B>0$.
There is $c=c(\chi,B)>0$ such that, uniformly for $X\geq3$,
$1\leq Y\leq X$, positive integers $q\leq(\log X)^B$, and Dirichlet
characters $\rho$ modulo $q$,
\begin{equation}\label{eq:gaussian-hecke-sw-estimate}
 \sum_{\mathrm N\mathfrak a\leq Y}
   \Lambda_E(\mathfrak a)\chi(\mathfrak a)\rho(\mathrm N\mathfrak a)
 =r_{\chi,\rho}Y+
 O_{\chi,B}\bigl(X\exp(-c\sqrt{\log X})\bigr).
\end{equation}
Here $\Lambda_E(\mathfrak p^j)=\log\mathrm N\mathfrak p$ for a prime
ideal $\mathfrak p$ and $j\geq1$, and it is zero otherwise; characters
are extended by zero on ideals not coprime to their moduli.
The number $r_{\chi,\rho}$ is one if the primitive Hecke character
inducing $\chi(\rho\circ\mathrm N)$ is trivial, and zero otherwise.
There is also $c_B>0$ such that, with the same ranges of $X,Y,q,\rho$,
\begin{equation}\label{eq:rational-sw-estimate}
 \sum_{n\leq Y}\Lambda(n)\rho(n)
 =1_{\rho\ \mathrm{principal}}Y+
 O_B\bigl(X\exp(-c_B\sqrt{\log X})\bigr).
\end{equation}
In particular, either error is $O(X(\log X)^{-D})$ for every fixed
$D>0$, with the corresponding dependence on $D$ allowed.
\end{lemma}

\begin{proof}
We use the published fixed-field estimate of Mitsui
\cite[\S1, Lemma~15, p.~12]{Mitsui1956}.
In his notation, for a fixed number field $K$, a fixed $A>0$, and any
ray character $\vartheta$ whose finite modulus has norm at most
$(\log x)^A$, it states
\begin{equation}\label{eq:mitsui-weighted-primes}
 \sum_{\mathrm N\mathfrak p\leq x}
       \vartheta(\mathfrak p)\log\mathrm N\mathfrak p
 =E(\vartheta)x+
 O_{K,A}\bigl(x\exp(-c_{K,A}\sqrt{\log x})\bigr),
\end{equation}
where $E(\vartheta)=1$ for the principal character and is zero otherwise.
The sum is over prime ideals of every residue degree.
Mitsui defines these characters and $E(\vartheta)$ on p.~2, includes
imprimitive characters, and states the uniformity in the character
and its modulus immediately before Lemma~14 on p.~11.
Infinite primes may occur in the ray modulus; only its finite part
enters the norm bound.

For completeness, the elementary passage from Mitsui's proved
smoothed estimate (Lemma~14, pp.~11--12) to
\eqref{eq:mitsui-weighted-primes} also applies to complex characters.
Put
\[
 P(x,\vartheta)=\sum_{\mathrm N\mathfrak p\leq x}
  \vartheta(\mathfrak p)\log\mathrm N\mathfrak p
                         \log\frac{x}{\mathrm N\mathfrak p}.
\]
Lemma~14 gives $P(x,\vartheta)=E(\vartheta)x+
O_{K,A}(xe^{-c_1\sqrt{\log x}})$.
Subtract its values at $xe^h$ and $x$, where
$h=e^{-c_1\sqrt{\log x}/2}$, and divide by $h$.
The resulting left side is the prime sum in
\eqref{eq:mitsui-weighted-primes} plus boundary terms of absolute
value at most
\[
 \sum_{x<\mathrm N\mathfrak p\leq xe^h}\log\mathrm N\mathfrak p
       \ll_K(xh+1)\log x.
\]
Indeed, for each integer norm there are at most $[K:\Q]$ prime ideals.
The main term is $E(\vartheta)x(e^h-1)/h=E(\vartheta)x+O(xh)$,
and the error divided by $h$ is
$O_{K,A}(xe^{-c_1\sqrt{\log x}/2})$.
Reducing the positive exponent constant proves
\eqref{eq:mitsui-weighted-primes}.

Prime powers of exponent at least two contribute
$O_K(\sqrt x\log^2(2x))$: sum over exponents $j\geq2$ and use the
elementary bound $O_K(x^{1/j})$ for the number of prime ideals of norm
at most $x^{1/j}$.
Consequently \eqref{eq:mitsui-weighted-primes} holds with the prime
sum replaced by
$\sum_{\mathrm N\mathfrak a\leq x}
\Lambda_K(\mathfrak a)\vartheta(\mathfrak a)$.

The character $\rho\circ\mathrm N$ on $E$ has ray modulus $(q)$:
if $\alpha\equiv1\pmod{q\Z[i]}$, then
$\mathrm N\alpha\equiv1\pmod q$.
Thus $\chi(\rho\circ\mathrm N)$ has a modulus dividing
$\mathfrak m(q)$, whose norm is $\mathrm N\mathfrak m\,q^2$.
For $\sqrt X\leq Y\leq X$, the bounds
$q\leq(\log X)^B$ and $\log Y\geq\frac12\log X$ imply
\[
 \mathrm N\mathfrak m\,q^2\leq(\log Y)^{2B+1}
\]
for all sufficiently large $X$, with the threshold depending only
on $\mathfrak m$ and $B$.
Apply the preceding estimate with $K=E$, $x=Y$, and $A=2B+1$.
Its error is at most the error in
\eqref{eq:gaussian-hecke-sw-estimate}, after reducing $c$.
For $1\leq Y<\sqrt X$, the elementary bound
$\sum_{\mathrm N\mathfrak a\leq Y}\Lambda_E(\mathfrak a)
\ll_E Y\log^2(2Y)$ suffices, and the remaining bounded values of
$X$ are absorbed in the implied constant.

The principal character at the chosen modulus is exactly the
character induced by the trivial primitive character, giving the
stated $r_{\chi,\rho}$.
Equivalently one can work at the primitive conductor and remove
the extra Euler factors: the corresponding prime-power towers
contribute only
$O_\chi(\log(q+2)\log(Y+2))$, which is absorbed in the error.

For \eqref{eq:rational-sw-estimate}, take $K=\Q$ and use the ray
character attached to $\rho$, allowing the real infinite place
for odd characters.  Its finite modulus norm is $q$.
The same proof uses $A=B+1$ for $Y\geq\sqrt X$ and the elementary
bound for smaller $Y$.  Finally,
$e^{-c\sqrt{\log X}}\ll_{c,D}(\log X)^{-D}$ for every fixed $D>0$.
\end{proof}
